\begin{enumerate}[label=\thesubsection.\arabic*, ref=\thesubsection.\theenumi]
\item Let $k \neq 0$. The binary operation $*$ on $\mathbb{R}-\cbrak{0}$ defined by
\begin{align}
a * b = \frac{ab}{k}
\label{eq:EE26BTECHXXXXX-star}
\end{align}
is associative, has $k$ as its identity, and with $\phi\brak{a} = \frac{a}{k}$,
\begin{align}
\phi\brak{a * b} = \phi\brak{a}\phi\brak{b}
\label{eq:EE26BTECHXXXXX-star-phi}
\end{align}
\solution
From \eqref{eq:EE26BTECHXXXXX-star},
\begin{align}
\phi\brak{a * b} = \frac{1}{k}\cdot\frac{ab}{k} = \frac{a}{k}\cdot\frac{b}{k}
\end{align}
yielding \eqref{eq:EE26BTECHXXXXX-star-phi}. Since $\phi$ is one-to-one,
\begin{align}
\phi\brak{\brak{a*b}*c} = \phi\brak{a}\phi\brak{b}\phi\brak{c} = \phi\brak{a*\brak{b*c}}
\end{align}
and $\phi\brak{a * k} = \phi\brak{a}\phi\brak{k} = \phi\brak{a}$, so $k$ is the identity.
\item For $u \in \sbrak{-1,1}$,
\begin{align}
\sin^{-1}u + \cos^{-1}u = \frac{\pi}{2}
\label{eq:EE26BTECHXXXXX-arcsin-comp}
\end{align}
\solution
Let $\theta = \sin^{-1}u \in \sbrak{-\frac{\pi}{2},\frac{\pi}{2}}$. Then $\frac{\pi}{2}-\theta \in \sbrak{0,\pi}$ and
\begin{align}
\cos\brak{\frac{\pi}{2}-\theta} = \sin\theta = u
\implies \cos^{-1}u = \frac{\pi}{2}-\theta
\end{align}
yielding \eqref{eq:EE26BTECHXXXXX-arcsin-comp}.
\item The principal value of $\cot^{-1}$ lies in $\brak{0,\pi}$ and
\begin{align}
\cot^{-1}\brak{\cot\theta} = \theta, \quad \theta \in \brak{0,\pi}
\label{eq:EE26BTECHXXXXX-pv-cot}
\end{align}
Similarly,
\begin{align}
\tan^{-1}\brak{\tan\theta} = \theta, \quad \theta \in \brak{-\frac{\pi}{2},\frac{\pi}{2}}
\end{align}
\item For any real $x$,
\begin{align}
1 \pm \sin x = \brak{\cos\frac{x}{2} \pm \sin\frac{x}{2}}^2
\label{eq:EE26BTECHXXXXX-half-angle}
\end{align}
\solution
Expanding the right hand side,
\begin{align}
\cos^2\frac{x}{2} + \sin^2\frac{x}{2} \pm 2\sin\frac{x}{2}\cos\frac{x}{2} = 1 \pm \sin x
\end{align}
\item Let $\phi: S \to \mathbb{R}^m$ be any function. The relation
\begin{align}
\vec{x} R \vec{y} \iff \phi\brak{\vec{x}} = \phi\brak{\vec{y}}
\label{eq:EE26BTECHXXXXX-kernel-rel}
\end{align}
is an equivalence relation and the equivalence class of $\vec{x}$ is
\begin{align}
\sbrak{\vec{x}} = \cbrak{\vec{y} \in S : \phi\brak{\vec{y}} = \phi\brak{\vec{x}}}
\end{align}
\solution
\begin{align}
\phi\brak{\vec{x}} &= \phi\brak{\vec{x}} && \text{(reflexive)}
\\
\phi\brak{\vec{x}} = \phi\brak{\vec{y}} &\implies \phi\brak{\vec{y}} = \phi\brak{\vec{x}} && \text{(symmetric)}
\\
\phi\brak{\vec{x}} = \phi\brak{\vec{y}},\, \phi\brak{\vec{y}} = \phi\brak{\vec{z}} &\implies \phi\brak{\vec{x}} = \phi\brak{\vec{z}} && \text{(transitive)}
\end{align}
\item For $n\times n$ matrices,
\begin{align}
\mydet{\vec{AB}} = \mydet{\vec{A}}\mydet{\vec{B}}
\label{eq:EE26BTECHXXXXX-det-product}
\end{align}
\solution
Any invertible $\vec{A}$ is a product $\vec{E}_1\vec{E}_2\cdots\vec{E}_k$ of elementary matrices and, for a single elementary matrix, it is easy to verify that $\mydet{\vec{E}\vec{B}} = \mydet{\vec{E}}\mydet{\vec{B}}$. Hence
\begin{align}
\mydet{\vec{AB}} = \mydet{\vec{E}_1}\cdots\mydet{\vec{E}_k}\mydet{\vec{B}}
= \mydet{\vec{A}}\mydet{\vec{B}}
\end{align}
upon substituting $\vec{B} = \vec{I}$. If $\vec{A}$ is singular, so is $\vec{AB}$ and both sides are $0$.
\item Adding a multiple of one row of a matrix to another row does not change its determinant.
\begin{align}
\mydet{\vec{E}\vec{A}} = \mydet{\vec{A}}, \quad \vec{E} = \vec{I} + \lambda\vec{e}_i\vec{e}_j^{\top}, \, i \neq j
\label{eq:EE26BTECHXXXXX-det-rowop}
\end{align}
Also, a common factor of a row can be taken out of the determinant, and the determinant can be expanded along any row.
\solution
$\vec{E}$ is triangular with unit diagonal, so $\mydet{\vec{E}} = 1$. Substituting in \eqref{eq:EE26BTECHXXXXX-det-product}, we get \eqref{eq:EE26BTECHXXXXX-det-rowop}.
\item The scalar triple product of three vectors is
\begin{align}
\sbrak{\vec{a}, \vec{b}, \vec{c}} = \vec{a}^{\top}\brak{\vec{b}\times\vec{c}} = \mydet{\vec{a} & \vec{b} & \vec{c}}
\label{eq:EE26BTECHXXXXX-triple}
\end{align}
\solution
Expanding $\mydet{\vec{a} & \vec{b} & \vec{c}}$ along the first column, the cofactors of $\vec{a}$ are the components of $\vec{b}\times\vec{c}$.
\item The scalar projection of $\vec{a}$ on $\vec{b}$ is
\begin{align}
\frac{\vec{a}^{\top}\vec{b}}{\norm{\vec{b}}}
\label{eq:EE26BTECHXXXXX-proj}
\end{align}
Also,
\begin{align}
\norm{\vec{a}+\vec{b}}^2 = \norm{\vec{a}}^2 + 2\vec{a}^{\top}\vec{b} + \norm{\vec{b}}^2,
\quad
\cos\theta = \frac{\vec{a}^{\top}\vec{b}}{\norm{\vec{a}}\norm{\vec{b}}}
\label{eq:EE26BTECHXXXXX-angle}
\end{align}
\solution
Write $\vec{a} = \lambda\vec{b} + \vec{a}_{\perp}$ with $\vec{b}^{\top}\vec{a}_{\perp} = 0$. Then
\begin{align}
\lambda = \frac{\vec{a}^{\top}\vec{b}}{\norm{\vec{b}}^2}
\implies \norm{\lambda\vec{b}} = \frac{\vec{a}^{\top}\vec{b}}{\norm{\vec{b}}}
\end{align}
The second result follows from expanding $\brak{\vec{a}+\vec{b}}^{\top}\brak{\vec{a}+\vec{b}}$.
\item A plane with normal vector $\vec{n}$ is
\begin{align}
\vec{n}^{\top}\vec{x} = c
\label{eq:EE26BTECHXXXXX-plane}
\end{align}
A plane parallel to \eqref{eq:EE26BTECHXXXXX-plane} and through $\vec{P}_0$ is $\vec{n}^{\top}\vec{x} = \vec{n}^{\top}\vec{P}_0$. The plane through $\vec{A}, \vec{B}, \vec{C}$ has a normal vector $\vec{n}$ satisfying
\begin{align}
\myvec{\brak{\vec{B}-\vec{A}}^{\top} \\ \brak{\vec{C}-\vec{A}}^{\top}}\vec{n} = \vec{0}, \quad c = \vec{n}^{\top}\vec{A}
\label{eq:EE26BTECHXXXXX-plane-3pts}
\end{align}
The distance of $\vec{P}$ from \eqref{eq:EE26BTECHXXXXX-plane} is
\begin{align}
d = \frac{\abs{\vec{n}^{\top}\vec{P} - c}}{\norm{\vec{n}}}
\label{eq:EE26BTECHXXXXX-plane-dist}
\end{align}
\solution
The foot of the perpendicular is $\vec{Q} = \vec{P} - \lambda\vec{n}$. Since $\vec{n}^{\top}\vec{Q} = c$,
\begin{align}
\vec{n}^{\top}\vec{P} - \lambda\norm{\vec{n}}^2 = c
\implies \lambda = \frac{\vec{n}^{\top}\vec{P} - c}{\norm{\vec{n}}^2}
\end{align}
yielding \eqref{eq:EE26BTECHXXXXX-plane-dist} from $d = \norm{\vec{P}-\vec{Q}} = \abs{\lambda}\norm{\vec{n}}$. The normal vector in \eqref{eq:EE26BTECHXXXXX-plane-3pts} is perpendicular to two directions in the plane.
\item The lines $\vec{x} = \vec{A}_1 + \lambda\vec{m}_1$ and $\vec{x} = \vec{A}_2 + \mu\vec{m}_2$ intersect if and only if
\begin{align}
\vec{M}\myvec{\lambda \\ \mu} = \vec{A}_2 - \vec{A}_1, \quad \vec{M} = \myvec{\vec{m}_1 & -\vec{m}_2}
\label{eq:EE26BTECHXXXXX-line-isect}
\end{align}
is consistent, i.e.
\begin{align}
\rank\brak{\vec{M}} = \rank\brak{\myvec{\vec{M} & \vec{A}_2 - \vec{A}_1}}
\label{eq:EE26BTECHXXXXX-consistent}
\end{align}
If $\rank\brak{\vec{M}} = 2$, the point of intersection is unique.
\solution
At a common point, $\vec{A}_1 + \lambda\vec{m}_1 = \vec{A}_2 + \mu\vec{m}_2$, which yields \eqref{eq:EE26BTECHXXXXX-line-isect}. This system has a solution if and only if $\vec{A}_2 - \vec{A}_1$ lies in the column space of $\vec{M}$, which is \eqref{eq:EE26BTECHXXXXX-consistent}. If the columns are independent, the solution is unique.
\item If $\mydet{\vec{A}} \neq 0$, then $\vec{Ax} = \vec{b}$ has the unique solution
\begin{align}
\vec{x} = \vec{A}^{-1}\vec{b}, \quad \vec{A}^{-1} = \frac{\text{adj}\brak{\vec{A}}}{\mydet{\vec{A}}}
\label{eq:EE26BTECHXXXXX-linsolve}
\end{align}
\solution
Multiplying $\vec{Ax} = \vec{b}$ on the left by $\vec{A}^{-1}$ yields \eqref{eq:EE26BTECHXXXXX-linsolve}. The inverse exists since $\mydet{\vec{A}} \neq 0$.
\item The conditional probability of $A$ given $B$ is
\begin{align}
\pr{\cond{A}{B}} = \frac{\pr{AB}}{\pr{B}}, \quad \pr{B} > 0
\label{eq:EE26BTECHXXXXX-cond}
\end{align}
If $E_1, E_2, \dots, E_n$ are mutually exclusive and exhaustive events, then (Bayes' theorem)
\begin{align}
\pr{\cond{E_i}{A}} = \frac{\pr{\cond{A}{E_i}}\pr{E_i}}{\sum_{j=1}^{n}\pr{\cond{A}{E_j}}\pr{E_j}}
\label{eq:EE26BTECHXXXXX-bayes}
\end{align}
\solution
From \eqref{eq:EE26BTECHXXXXX-cond}, $\pr{E_iA} = \pr{\cond{A}{E_i}}\pr{E_i}$. Since $A = \sum_j E_jA$ and the $E_jA$ are mutually exclusive, from \eqref{eq:EE26BTECHXXXXX-Prob_3},
\begin{align}
\pr{A} = \sum_{j=1}^{n}\pr{\cond{A}{E_j}}\pr{E_j}
\end{align}
Substituting in $\pr{\cond{E_i}{A}} = \frac{\pr{E_iA}}{\pr{A}}$ yields \eqref{eq:EE26BTECHXXXXX-bayes}.
\item For mutually exclusive events $A$ and $B$,
\begin{align}
\pr{A + B} = \pr{A} + \pr{B} \label{eq:EE26BTECHXXXXX-Prob_3}
\end{align}
\item For $y > 0$,
\begin{align}
\frac{y^{\prime}}{y} = \brak{\ln y}^{\prime}, \quad
\frac{y^{\prime\prime}}{y} - \brak{\frac{y^{\prime}}{y}}^2 = \brak{\ln y}^{\prime\prime}
\label{eq:EE26BTECHXXXXX-logder}
\end{align}
\solution
Using the quotient rule,
\begin{align}
\brak{\frac{y^{\prime}}{y}}^{\prime} = \frac{y^{\prime\prime}y - \brak{y^{\prime}}^2}{y^2}
= \frac{y^{\prime\prime}}{y} - \brak{\frac{y^{\prime}}{y}}^2
\end{align}
\item If $f^{\prime}\brak{x} > 0$ on an interval $I$, $f$ is strictly increasing on $I$. If $f^{\prime}\brak{x} < 0$ on $I$, $f$ is strictly decreasing on $I$. Also, if $f^{\prime}\brak{c} = 0$ and $f^{\prime\prime}\brak{c} < 0$, $f$ has a maximum at $c$.
\label{eq:EE26BTECHXXXXX-monotone}
\\
\solution
For $a < b$ in $I$, by the mean value theorem there is a $\xi \in \brak{a,b}$ such that
\begin{align}
f\brak{b} - f\brak{a} = f^{\prime}\brak{\xi}\brak{b-a}
\end{align}
whose sign is the sign of $f^{\prime}\brak{\xi}$.
\item The Laplace transform of $f\brak{x}$, $x \ge 0$, is
\begin{align}
F\brak{s} = \mathcal{L}\cbrak{f\brak{x}} = \int_{0}^{\infty} f\brak{x}e^{-sx}\, dx
\label{eq:EE26BTECHXXXXX-laplace-def}
\end{align}
It is linear, one-to-one, and
\begin{align}
\mathcal{L}\cbrak{f^{\prime}\brak{x}} &= sF\brak{s} - f\brak{0}
\label{eq:EE26BTECHXXXXX-laplace-deriv}
\\
\mathcal{L}\cbrak{xf\brak{x}} &= -F^{\prime}\brak{s}
\label{eq:EE26BTECHXXXXX-laplace-xmult}
\end{align}
\solution
Integrating by parts,
\begin{align}
\int_{0}^{\infty} f^{\prime}e^{-sx}\, dx = \sbrak{fe^{-sx}}_{0}^{\infty} + s\int_{0}^{\infty} fe^{-sx}\, dx = -f\brak{0} + sF\brak{s}
\end{align}
yielding \eqref{eq:EE26BTECHXXXXX-laplace-deriv}. Differentiating \eqref{eq:EE26BTECHXXXXX-laplace-def} with respect to $s$,
\begin{align}
F^{\prime}\brak{s} = \int_{0}^{\infty} \brak{-x}f\brak{x}e^{-sx}\, dx
\end{align}
which is \eqref{eq:EE26BTECHXXXXX-laplace-xmult}.
\item Standard transforms, for $s > a$,
\begin{align}
\mathcal{L}\cbrak{1} = \frac{1}{s}, \quad
\mathcal{L}\cbrak{e^{ax}} = \frac{1}{s-a}, \quad
\mathcal{L}\cbrak{xe^{ax}} = \frac{1}{\brak{s-a}^2}
\label{eq:EE26BTECHXXXXX-laplace-table}
\\
\mathcal{L}\cbrak{\ln x} = -\frac{\gamma + \ln s}{s}
\label{eq:EE26BTECHXXXXX-laplace-ln}
\end{align}
where $\gamma$ is Euler's constant.
\solution
Since $\int_{0}^{\infty} e^{-\brak{s-a}x}\, dx = \frac{1}{s-a}$, the first two follow. From \eqref{eq:EE26BTECHXXXXX-laplace-xmult},
\begin{align}
\mathcal{L}\cbrak{xe^{ax}} = -\frac{d}{ds}\brak{\frac{1}{s-a}} = \frac{1}{\brak{s-a}^2}
\end{align}
For $b > -1$, substituting $u = sx$,
\begin{align}
\mathcal{L}\cbrak{x^b} = \int_{0}^{\infty} x^be^{-sx}\, dx = \frac{1}{s^{b+1}}\int_{0}^{\infty} u^be^{-u}\, du = \frac{\Gamma\brak{b+1}}{s^{b+1}}
\end{align}
Differentiating with respect to $b$ at $b = 0$, using $\Gamma\brak{1} = 1$ and $\Gamma^{\prime}\brak{1} = -\gamma$,
\begin{align}
\mathcal{L}\cbrak{\ln x} = \frac{\Gamma^{\prime}\brak{1} - \Gamma\brak{1}\ln s}{s}
\end{align}
yielding \eqref{eq:EE26BTECHXXXXX-laplace-ln}.
\item For $\abs{x} < 1$, the substitution $x = \tanh t$ gives
\begin{align}
\frac{dx}{dt} = 1 - x^2, \quad t = \frac{1}{2}\ln\frac{1+x}{1-x}
\label{eq:EE26BTECHXXXXX-tanh-sub}
\end{align}
\solution
Since $\frac{d}{dt}\tanh t = 1 - \tanh^2 t$, the first result follows. Solving $x = \frac{e^{2t}-1}{e^{2t}+1}$ for $e^{2t}$ gives $e^{2t} = \frac{1+x}{1-x}$.
\item Standard integrals.
\begin{align}
\int x^n\, dx &= \frac{x^{n+1}}{n+1} + C, \quad n \neq -1
\\
\int \frac{dx}{a^2 + x^2} &= \frac{1}{a}\tan^{-1}\frac{x}{a} + C, \quad a \neq 0
\label{eq:EE26BTECHXXXXX-stdint}
\\
\int \sec^2 x\, dx &= \tan x + C, \quad \int \cosec^2 x\, dx = -\cot x + C
\end{align}
\solution
For \eqref{eq:EE26BTECHXXXXX-stdint}, substituting $x = a\tan\theta$, $dx = a\sec^2\theta\, d\theta$,
\begin{align}
\int \frac{a\sec^2\theta}{a^2\brak{1+\tan^2\theta}}\, d\theta = \frac{1}{a}\int d\theta = \frac{\theta}{a} + C
\end{align}
\item For any $a, b$,
\begin{align}
a^3 + b^3 = \brak{a+b}^3 - 3ab\brak{a+b}
\label{eq:EE26BTECHXXXXX-cube-sum}
\end{align}
\solution
Expanding $\brak{a+b}^3 = a^3 + 3a^2b + 3ab^2 + b^3$ yields \eqref{eq:EE26BTECHXXXXX-cube-sum}.
\item (King's property)
\begin{align}
\int_{a}^{b} f\brak{x}\, dx = \int_{a}^{b} f\brak{a+b-x}\, dx
\label{eq:EE26BTECHXXXXX-king}
\end{align}
\solution
Substituting $t = a + b - x$, $dx = -dt$, the limits $x = a, b$ become $t = b, a$.
\begin{align}
\int_{a}^{b} f\brak{x}\, dx = -\int_{b}^{a} f\brak{t}\, dt = \int_{a}^{b} f\brak{t}\, dt
\end{align}
\item In polar coordinates $x = r\cos\theta$, $y = r\sin\theta$,
\begin{align}
\iint_{D} dx\, dy = \iint_{D} r\, dr\, d\theta
\label{eq:EE26BTECHXXXXX-polar-area}
\end{align}
\solution
The Jacobian determinant is
\begin{align}
\mydet{\frac{\partial x}{\partial r} & \frac{\partial x}{\partial \theta} \\ \frac{\partial y}{\partial r} & \frac{\partial y}{\partial \theta}}
= \mydet{\cos\theta & -r\sin\theta \\ \sin\theta & r\cos\theta} = r
\end{align}
\item (Corner point theorem) The maximum of $Z = \vec{c}^{\top}\vec{x}$ over a bounded polygon with vertices $\vec{v}_1, \dots, \vec{v}_m$ is attained at a vertex.
\label{eq:EE26BTECHXXXXX-corner}
\\
\solution
Any point of the polygon is $\vec{x} = \sum_i \theta_i\vec{v}_i$, with $\theta_i \ge 0$, $\sum_i\theta_i = 1$. Then
\begin{align}
\vec{c}^{\top}\vec{x} = \sum_i \theta_i\vec{c}^{\top}\vec{v}_i \le \max_i \vec{c}^{\top}\vec{v}_i
\end{align}
\item For a right circular cone of slant height $l$, semi-vertical angle $\alpha$, radius $r$ and height $h$,
\begin{align}
r = l\sin\alpha, \quad h = l\cos\alpha, \quad
V = \frac{\pi}{3}r^2h = \frac{\pi}{3}\brak{l^2 - h^2}h
\label{eq:EE26BTECHXXXXX-cone}
\end{align}
\solution
The axis, radius and slant edge form a right triangle with hypotenuse $l$, so $r^2 + h^2 = l^2$, from which the above follows.
\end{enumerate}
